% TEMPLATE-MILC-v3.tex - plantilla maestra UNICA de las guias MILC v3.
%
% La usan las tres familias de guias, cada una con su driver:
%   · Grados 6-11  -> scripts/build-guias-g11.py
%   · Bebras       -> scripts/build-guias-bebras.py
%   · Semillero    -> scripts/build-guias-semillero.py
%
% Antes habia tres copias casi identicas (~400 lineas iguales, 6 distintas):
% cada mejora habia que aplicarla tres veces, y en la practica no se hacia
% (el motor de recursos vivia solo en dos de ellas). Lo que varia entre
% familias esta parametrizado en 6 placeholders que cada driver llena
% (se nombran aqui SIN su sintaxis real: el builder reemplaza por texto plano,
% tambien dentro de los comentarios, y un valor multilinea rompe el preambulo):
%
%   HEADER_IZQ        encabezado izquierdo de pagina
%   HEADER_DER        encabezado derecho de pagina
%   PORTADA_ETIQUETA  rotulo sobre el numero grande (GRADO/MOMENTO/MODULO)
%   PORTADA_SERIE     linea de serie de la portada
%   PORTADA_ID        identificador de la guia en la portada
%   PORTADA_CIRCULO   el circulito de grado (°), o vacio si no aplica
%
% El resto de claves las documenta content/guias/_SCHEMA.md. El script valida
% que no queden placeholders sin reemplazar antes de escribir el archivo final.

\documentclass[11pt,letterpaper]{article}
\usepackage[margin=2.0cm]{geometry}
\usepackage[table]{xcolor}
\usepackage{fontspec}
\usepackage[spanish]{babel}
\usepackage{tikz}
% Librerías TikZ para los diagramas de las guías (bloque `recursos:`):
%  · babel  → babel en español vuelve activos caracteres como «>», y TikZ los usa
%             en sus opciones (>=stealth, ->). Sin esta librería, cualquier
%             diagrama con flechas revienta con «\language@active@arg> has an
%             extra }». Debe ir ANTES de que se dibuje nada.
%  · positioning        → sintaxis «below left=2pt and 3pt».
%  · decorations.pathreplacing → llaves ({brace}) para señalar rangos.
\usetikzlibrary{babel,positioning,decorations.pathreplacing}
\usepackage{graphicx}
% Circuitos: el trabajo de electrónica se hace hoy con micro:bit y sus sensores
% integrados (MakeCode, sin cableado), así que no se necesitan esquemáticos.
% Si algún día entran montajes con protoboard, basta descomentar esta línea:
% los diagramas .tex/.tikz declarados en `recursos:` ya se incrustan solos.
% \usepackage{circuitikz}
\usepackage[most]{tcolorbox}
\usepackage{tabularx}
\usepackage{array}
\usepackage{fancyhdr}
\usepackage{titlesec}
\usepackage[document]{ragged2e}  % bandera a la derecha en todo el cuerpo
\usepackage{enumitem}
\usepackage{microtype}

% Helvetica Neue no trae la flecha U+2192, asi que cada «→» escrito literalmente
% en el YAML se compilaba a NADA, en silencio (462 flechas en 61 guias: rutas del
% tipo «paleta → tipografia → lockup» salian rotas). Mapeamos el caracter a su
% version matematica, que si tiene glifo. Asi el autor puede seguir escribiendo
% «→» comodamente en el YAML.
\usepackage{newunicodechar}
\newunicodechar{→}{\ensuremath{\rightarrow}}
% Mismo problema con los simbolos de marca y vinetas: el «✓» de «una palomita ✓»
% salia como cuadrito vacio en el PDF de todas las guias que lo usaban.
\usepackage{pifont}
\newunicodechar{✓}{\ding{51}}
\newunicodechar{✗}{\ding{55}}
\newunicodechar{✔}{\ding{52}}
\newunicodechar{•}{\textbullet}
\newunicodechar{≥}{\ensuremath{\geq}}
\newunicodechar{≤}{\ensuremath{\leq}}

\setmainfont{Helvetica Neue}
\setsansfont{Helvetica Neue}
\setlength{\parindent}{0pt}
\setlength{\parskip}{6pt}
% Sin particion de palabras: en 6.º-7.º una palabra cortada por guion al final
% de linea («Comuni-cacion») frena la lectura mucho mas que una linea corta.
\hyphenpenalty=10000
\exhyphenpenalty=10000
\pretolerance=2000
\tolerance=3000
\setlength{\emergencystretch}{3em}
\linespread{1.10}
\renewcommand{\arraystretch}{1.25}
\setlist{leftmargin=5mm,itemsep=1mm,topsep=1mm}
\newcolumntype{Y}{>{\RaggedRight\arraybackslash}X}

\definecolor{milcMagenta}{HTML}{E6007E}
\definecolor{milcVino}{HTML}{5A0038}
\definecolor{milcTurquesa}{HTML}{0093A5}
\definecolor{milcVerde}{HTML}{58A923}
\definecolor{milcMostaza}{HTML}{E5B400}
\definecolor{milcOcre}{HTML}{B86600}
\definecolor{milcNegro}{HTML}{241816}
\definecolor{milcGris}{HTML}{F7F7F4}
\definecolor{milcLinea}{HTML}{E8E2DD}
\definecolor{milcGradePrimary}{HTML}{0066FF}
\definecolor{milcGradeDark}{HTML}{003D99}
\definecolor{milcGradeSoft}{HTML}{D6E8FF}
\definecolor{milcGradeAccent}{HTML}{CFE7FF}
\definecolor{milcGradeText}{HTML}{FFFFFF}
\color{milcNegro}

\pagestyle{fancy}
\fancyhf{}
\lhead{\small Tecnología e Informática · Grado 6}
\rhead{\small Guía 16 · MILC v3}
\cfoot{\small\thepage}
\renewcommand{\headrulewidth}{0pt}

\newtcolorbox{titlebox}[1]{
  enhanced,colback=#1,colframe=#1,arc=7mm,boxrule=0pt,
  left=7mm,right=7mm,top=3mm,bottom=3mm,
  before skip=8pt,after skip=8pt,
  fontupper=\sffamily\bfseries\Large\color{white}
}

\newtcolorbox{softbox}[3]{
  enhanced,colback=#2,colframe=#1,borderline west={4pt}{0pt}{#1},
  arc=2mm,boxrule=.4pt,left=4mm,right=4mm,top=3mm,bottom=3mm,
  before skip=6pt,after skip=8pt,title={#3},coltitle=white,
  colbacktitle=#1,fonttitle=\sffamily\bfseries,fontupper=\RaggedRight,
  attach boxed title to top left={xshift=3mm,yshift=-2mm},
  boxed title style={arc=2mm, boxrule=0pt}
}

\newtcolorbox{drawbox}[2]{
  enhanced,colback=white,colframe=#1,arc=4mm,boxrule=1pt,
  left=5mm,right=5mm,top=4mm,bottom=4mm,before skip=8pt,after skip=8pt,
  title={#2},coltitle=white,colbacktitle=#1,fonttitle=\sffamily\bfseries,
  fontupper=\RaggedRight,
  attach boxed title to top left={xshift=4mm,yshift=-2mm},
  boxed title style={arc=3mm, boxrule=0pt}
}

\newtcolorbox{infoband}[1]{
  enhanced,colback=#1!8,colframe=#1!50,arc=2mm,boxrule=.5pt,
  left=4mm,right=4mm,top=2mm,bottom=2mm,before skip=4pt,after skip=4pt,
  fontupper=\small\RaggedRight
}

% Puente narrativo entre secciones (contrato editorial Conéctate).
% Da continuidad: "Acabas de X. Ahora vas a Y porque Z."
% Visualmente: banda izquierda en color del grado, texto en itálica pequeña.
\newtcolorbox{puentebox}{
  enhanced,colback=milcGradeSoft,colframe=milcGradeDark,
  borderline west={3pt}{0pt}{milcGradeDark},
  arc=0mm,boxrule=0pt,left=5mm,right=4mm,top=2.5mm,bottom=2.5mm,
  before skip=4pt,after skip=8pt,
  fontupper=\itshape\small\color{milcNegro}\RaggedRight
}
% La flecha va en modo matematico a proposito: Helvetica Neue NO tiene el glifo
% U+2192, asi que \textrightarrow se compilaba a NADA («Missing character: There
% is no → in font Helvetica Neue») y el puente salia sin flecha. En modo
% matematico la flecha viene de la fuente matematica y si se dibuja.
\newcommand{\puente}[1]{%
  \begin{puentebox}%
    \textcolor{milcGradeDark}{$\rightarrow$}\hspace{2mm}#1%
  \end{puentebox}%
}

\newcommand{\linead}[1]{\noindent\color{milcLinea}\rule{#1}{0.5pt}\color{milcNegro}}
\newcommand{\respuesta}[1]{\par\vspace{2mm}\foreach \n in {1,...,#1}{\noindent\color{milcLinea}\rule{\linewidth}{0.5pt}\par\vspace{5mm}}\color{milcNegro}}
\newcommand{\checkbox}{\textcolor{milcGradeDark}{\fbox{\phantom{x}}}\hspace{5pt}}

\newcommand{\phasecard}[4]{
\begin{tcolorbox}[enhanced,colback=#3,colframe=#2,arc=3mm,boxrule=.5pt,height=3.05cm,valign=center,halign=center,left=1.5mm,right=1.5mm,top=2mm,bottom=2mm]
{\sffamily\bfseries\fontsize{22}{24}\selectfont\textcolor{#2}{#1}}\par
{\sffamily\bfseries\small #4}
\end{tcolorbox}
}

% ── Piezas del rediseno 2026 (legibilidad 6.º-11.º) ───────────────────
% Banda de estacion: reemplaza el titlebox macizo. Titulo a la izquierda,
% dato practico (tiempo/modalidad) a la derecha, en una sola linea.
\newcommand{\milcestacion}[2]{%
\begin{tcolorbox}[enhanced,colback=#1,colframe=#1,arc=2mm,boxrule=0pt,
  left=5mm,right=5mm,top=2.6mm,bottom=2.6mm,before skip=11pt,after skip=7pt]
{\sffamily\bfseries\fontsize{15}{18}\selectfont\color{white}#2}
\end{tcolorbox}}

% Tarjeta de paso numerada: sustituye las tablas «Pilar / Aplicacion», que
% eran formato de consulta docente, no de estudiante. El numero va en un
% circulo sobre el borde para que la secuencia se lea de un vistazo.
\newcommand{\milcpaso}[3]{%
\begin{tcolorbox}[enhanced,colback=white,colframe=milcLinea,arc=1.5mm,boxrule=.9pt,
  left=14mm,right=4mm,top=3mm,bottom=3mm,before skip=4pt,after skip=4pt,
  fontupper=\RaggedRight,
  overlay={\node[anchor=north west,circle,fill=#2,text=white,inner sep=0pt,
    minimum size=7.4mm,font=\sffamily\bfseries\small]
    at ([xshift=3.6mm,yshift=-3.2mm]frame.north west) {#1};}]
#3
\end{tcolorbox}}

% Casilla marcable de tamano util con lapiz (la anterior era \fbox{\phantom{x}},
% de alto variable segun el texto que la seguia).
\newcommand{\milcmarca}{\framebox[3.8mm]{\rule{0pt}{3.4mm}}}

% Fila marcable de lista suelta (checks de la estacion 1).
\newcommand{\milccheckrow}[1]{%
\par\vspace{1.8mm}\noindent
\makebox[6.4mm][l]{\raisebox{-.55mm}{\textcolor{milcNegro}{\milcmarca}}}%
\parbox[t]{\dimexpr\linewidth-6.4mm\relax}{\RaggedRight #1}\par}

% Celda marcable dentro de la rubrica (tabla de autoevaluacion). Misma
% sangria francesa, pero pensada para vivir dentro de una columna X.
\newcommand{\milcopcion}[1]{%
\makebox[5.2mm][l]{\raisebox{-.55mm}{\textcolor{milcNegro}{\milcmarca}}}%
\parbox[t]{\dimexpr\linewidth-5.2mm\relax}{\RaggedRight #1}}

% ── Recursos visuales de la guía ──────────────────────────────────────
% Declarados en el bloque `recursos:` del YAML y emitidos por el builder.
% \guiaFigura   → imagen rasterizada o PDF (foto, carta celeste, montaje).
% \guiaDiagrama → fuente TikZ/circuitikz (.tex) incrustada como vector:
%                 es la vía para circuitos y esquemáticos editables.
% #1 = ruta relativa al .tex · #2 = pie (puede ir vacío)
\newcommand{\guiaFigura}[2]{%
\begin{center}
\includegraphics[width=0.88\linewidth,height=0.38\textheight,keepaspectratio]{#1}\\[1.5mm]
{\footnotesize\itshape\textcolor{milcNegro!75}{#2}}
\end{center}
}

\newcommand{\guiaDiagrama}[2]{%
\begin{center}
\input{#1}\\[1.5mm]
{\footnotesize\itshape\textcolor{milcNegro!75}{#2}}
\end{center}
}

\begin{document}
\thispagestyle{empty}

% ── Portada ───────────────────────────────────────────────────────────
% Antes: un rectangulo de color a pagina completa. Se veia bien en pantalla,
% pero en la fotocopiadora del colegio se comia el toner y salia gris sucio.
% Ahora la banda ocupa el tercio superior y el resto va en blanco.
\begin{tikzpicture}[remember picture,overlay]
  \path[fill=milcGradePrimary]
    (current page.north west) rectangle ([yshift=-10.6cm]current page.north east);

  \node[anchor=north west,text=milcGradeText] at ([xshift=2.0cm,yshift=-1.55cm]current page.north west)
    {\sffamily\bfseries\fontsize{12}{14}\selectfont GRADO};
  \node[anchor=north west,text=milcGradeText,opacity=.85] at ([xshift=2.0cm,yshift=-2.35cm]current page.north west)
    {\sffamily\bfseries\fontsize{8.5}{10}\selectfont SERIE GUÍAS · TECNOLOGÍA E INFORMÁTICA};
  \node[anchor=north east,text=milcGradeText,opacity=.85,align=right] at ([xshift=-2.0cm,yshift=-1.55cm]current page.north east)
    {\sffamily\bfseries\fontsize{9}{12}\selectfont I.E. Sor María Juliana\\Cartago, Valle del Cauca};

  \node[anchor=west,text=milcGradeText] (gradonum) at ([xshift=1.85cm,yshift=-7.1cm]current page.north west)
    {\sffamily\bfseries\fontsize{112}{112}\selectfont 6};
  % El circulito de grado (°) solo aplica a las guias de grado («11°»); en
  % Bebras y Semillero el numero grande es un momento/modulo, no un grado.
  % Se posiciona relativo al nodo `gradonum`, no en coordenadas absolutas.
  \draw[milcGradeText,line width=1.6pt]
    ([xshift=2.0mm,yshift=-4.5mm]gradonum.north east) circle (.15cm);

  \node[anchor=west,text=milcGradeText,text width=11.4cm,align=left] at ([xshift=1.5cm,yshift=0.5cm]gradonum.east)
    {
     {\sffamily\bfseries\fontsize{9}{11}\selectfont\MakeUppercase{Período 2 · Hardware y software}}\\[3mm]
     {\sffamily\bfseries\fontsize{21}{25}\selectfont\setlength{\baselineskip}{27pt}El sistema\\[4pt]operativo ---\\[4pt]el gerente\\[4pt]que coordina\\[4pt]todo}\\[3.5mm]
     {\sffamily\bfseries\fontsize{15}{18}\selectfont Guía 16-6-TIC}
    };
\end{tikzpicture}

\vspace*{9.4cm}
\vfill

{\sffamily\bfseries\fontsize{8}{10}\selectfont\textcolor{milcGradeDark}{LO QUE TE LLEVAS HOY}}

\begin{tcolorbox}[enhanced,colback=white,colframe=milcNegro,arc=2mm,boxrule=1.6pt,
  left=5mm,right=5mm,top=4mm,bottom=4mm,before skip=3pt,after skip=12pt,
  fontupper=\RaggedRight\fontsize{11}{15.5}\selectfont]
Un \textbf{tour anotado del sistema operativo} en una hoja del cuaderno con: \textbf{(a) dibujo de un escritorio típico de Windows} con 5 elementos clave etiquetados (escritorio, ventana abierta, barra de tareas, menú inicio, papelera), \textbf{(b) tabla de 5 funciones del SO} (coordinar hardware, gestionar archivos, ejecutar programas, manejar usuarios, controlar memoria), y \textbf{(c) tu opinión escrita} de cuál es el SO que más usas y por qué (Windows / Linux / Android / iOS / ChromeOS).
\end{tcolorbox}

\begin{tcolorbox}[enhanced,colback=milcGris,colframe=milcGris,arc=2mm,boxrule=0pt,
  left=5mm,right=5mm,top=3.5mm,bottom=3.5mm,after skip=12pt]
{\sffamily\bfseries\fontsize{8}{10}\selectfont\textcolor{milcNegro!55}{ESTA GUÍA ES DE}}\\[3mm]
\begin{tabularx}{\linewidth}{X p{3.2cm} p{3.2cm}}
\linead{\linewidth} & \linead{3.0cm} & \linead{3.0cm}\\[-1mm]
{\fontsize{8.5}{10}\selectfont\textcolor{milcNegro!55}{Nombre completo}} &
{\fontsize{8.5}{10}\selectfont\textcolor{milcNegro!55}{Grupo}} &
{\fontsize{8.5}{10}\selectfont\textcolor{milcNegro!55}{Fecha}}\\
\end{tabularx}
\end{tcolorbox}

\vfill

\noindent\textcolor{milcLinea}{\rule{\linewidth}{1pt}}\\[2mm]
{\sffamily\bfseries\fontsize{9.5}{12}\selectfont Autor: Dr. Álvaro Cárdenas Orozco}\\[1mm]
{\sffamily\fontsize{8.5}{11}\selectfont\textcolor{milcNegro!55}{Metodología MILC · apertura ancestral · pensamiento computacional · triángulo Dussel–estoicismo–Floridi}}

\newpage

\section*{Apertura: El sistema operativo --- el gerente que coordina toda la máquina}

\begin{softbox}{milcGradeDark}{milcGradeSoft}{Saber ancestral}
\textbf{En el taller grande de don Aurelio había varios aprendices, y todos sabían a quién obedecer.} En el centro de Cartago, en los años 60, hubo un taller de relojes donde trabajaban 4 aprendices al tiempo: uno limpiaba piezas, otro lijaba ejes, otro ensamblaba mecanismos, otro probaba el tic-tac final. \textbf{El que coordinaba todo era don Aurelio}, el maestro. Don Aurelio no hacía el trabajo manual; él \emph{decidía quién hacía qué, en qué momento, con qué herramienta}. Si llegaba un reloj urgente, asignaba a 2 aprendices a la vez. Si una pieza llegaba tarde, redistribuía las tareas. \textbf{Sin don Aurelio, los 4 aprendices se chocaban: dos limpiaban la misma pieza, otro buscaba la herramienta que su compañero tenía, nadie terminaba}. Con don Aurelio coordinando, el taller producía 10 relojes al día. Hoy ese coordinador en el computador se llama \textbf{sistema operativo}. Sin él, el hardware no se entiende con los programas, los archivos no se encuentran, los periféricos no obedecen. Es el gerente invisible que hace que todo funcione.
\end{softbox}

\begin{softbox}{milcTurquesa}{milcGris}{Saber contemporáneo}
El \textbf{sistema operativo (SO)} es el software más importante del computador. Es el \emph{primer programa} que se instala (después no se puede usar el equipo sin él) y el \emph{primero que enciende} cuando prendes el computador. Su trabajo es \textbf{coordinar todo lo demás}: hace que el teclado entienda con Word, que Chrome guarde tus descargas en la carpeta correcta, que Roblox use la memoria RAM sin pelearse con WhatsApp. Hay varios SO en el mundo. En computadores: \textbf{Windows} (el más común, de Microsoft), \textbf{macOS} (de Apple), \textbf{Linux} (gratuito, abierto), \textbf{ChromeOS} (de Google). En celulares: \textbf{Android} (Google), \textbf{iOS} (Apple). Aunque parecen distintos, \textbf{todos hacen lo mismo}: coordinan hardware, programas, archivos y usuarios. Aprender qué hace un SO te quita el miedo a probar uno nuevo --- son variantes del mismo gerente.
\end{softbox}

\begin{softbox}{milcMostaza}{milcGris}{Pregunta-puente}
\textit{Cuando prendes un computador, antes de que aparezca el escritorio, ves una pantalla que dice \emph{``Iniciando Windows''} (o ``Cargando macOS'', o el logo del fabricante). \textbf{¿Qué está haciendo el computador en esos segundos}? ¿Por qué no enciende y muestra todo de una vez?}
\end{softbox}

\begin{softbox}{milcVerde}{milcGris}{Saber y saber hacer}
Al terminar la clase: (1) podrás \textbf{identificar} 5 elementos típicos de un sistema operativo de escritorio; (2) sabrás \textbf{explicar} las 5 funciones que hace un SO; (3) podrás \textbf{distinguir} entre los principales SO del mercado (Windows, Linux, Android, iOS, ChromeOS); (4) habrás \textbf{aplicado} todo en un tour anotado de Windows con tu opinión personal.
\end{softbox}



\small
\begin{tabularx}{\textwidth}{>{\bfseries}p{3.0cm}Y}
\rowcolor{milcGradePrimary!12}Autor & Dr. Álvaro Cárdenas Orozco\\
DBA & Identifica componentes y sistemas tecnológicos en su entorno (MEN, T\&I 6°)\\
Referentes & Saber ancestral del relojero · Sistemas como cadena de oficios · Diagnóstico paso a paso\\
Producto final & Un \textbf{tour anotado del sistema operativo} en una hoja del cuaderno con: \textbf{(a) dibujo de un escritorio típico de Windows} con 5 elementos clave etiquetados (escritorio, ventana abierta, barra de tareas, menú inicio, papelera), \textbf{(b) tabla de 5 funciones del SO} (coordinar hardware, gestionar archivos, ejecutar programas, manejar usuarios, controlar memoria), y \textbf{(c) tu opinión escrita} de cuál es el SO que más usas y por qué (Windows / Linux / Android / iOS / ChromeOS).\\
\end{tabularx}
\normalsize

\puente{Hoy haces 4 cosas: reconoces los 5 elementos del escritorio, aprendes las 5 funciones del SO, recorres un computador real (o lo dibujas), y firmas tu trabajo con opinión.}

\milcestacion{milcGradeDark}{Ruta MILC: así vamos a aprender}
\begin{tabularx}{\textwidth}{>{\centering\arraybackslash}X>{\centering\arraybackslash}X>{\centering\arraybackslash}X>{\centering\arraybackslash}X}
\phasecard{1}{milcVerde}{milcGris}{Escucha\\[-1mm]\small Observo, escucho y pregunto.} &
\phasecard{2}{milcTurquesa}{milcGris}{Entender\\[-1mm]\small Sistematizo y ordeno.} &
\phasecard{3}{milcMagenta}{milcGris}{Hacer\\[-1mm]\small Construyo y aplico.} &
\phasecard{4}{milcVino}{milcGris}{Cierre\\[-1mm]\small Evalúo y reflexiono.}
\end{tabularx}


\puente{Empezamos identificando 5 elementos del escritorio de Windows que ya has visto antes.}

\milcestacion{milcVerde}{Estación 1 de 3 · Escucha}

\begin{softbox}{milcVerde}{milcGris}{Escena para escuchar}
\textbf{Actividad 1 · IDENTIFICA --- Reconoce 5 elementos del escritorio} (10 min · individual). Imagina que prendes un computador con Windows y aparece la pantalla principal. En el cuaderno, escribe el \textbf{nombre de 5 cosas} que aparecen ahí y describe cada una en 1 línea con tus palabras. Pista: hay 1 zona grande de fondo (\emph{el escritorio}), una barra abajo (\emph{barra de tareas}), un botón a la izquierda abajo (\emph{menú inicio}), un ícono de papelera, e \emph{íconos de programas o archivos}. Si no sabes el nombre técnico, describe lo que ves. Esta actividad recoge tu conocimiento previo --- no se vale ver la siguiente.
\end{softbox}

{\sffamily\bfseries\fontsize{8}{10}\selectfont\textcolor{milcNegro!55}{MARCA CUANDO LO HAYAS HECHO}}
\milccheckrow{Imaginé un escritorio de Windows.}
\milccheckrow{Listé 5 elementos diferentes con descripción propia.}
\milccheckrow{No miré la siguiente actividad ni busqué imágenes.}

\begin{infoband}{milcVerde}


\textbf{Lo que va al cuaderno (Actividad 1 · IDENTIFICA).} Título: «Actividad 1 --- 5 elementos del escritorio». \textbf{Formato:} lista numerada del 1 al 5. \textbf{Extensión:} 5 líneas. \textbf{Sabes que terminaste cuando} tienes 5 cosas diferentes nombradas y descritas con tus palabras.
\end{infoband}

\puente{Después del reconocimiento, aprendes qué hace el SO (las 5 funciones del gerente).}

\milcestacion{milcTurquesa}{Estación 2 de 3 · Entender}

\begin{softbox}{milcTurquesa}{milcGris}{Lo que vas a entender}
Lo que acabaste de listar son \textbf{los 5 elementos principales} del escritorio de un sistema operativo. Casi todos los SO de escritorio (Windows, macOS, Linux) tienen los mismos 5 elementos --- cambia el estilo, no la idea. Ahora vas a ponerles nombre técnico y entender qué hace cada uno.
\end{softbox}

\milcpaso{1}{milcTurquesa}{\textbf{Elemento 1 --- El escritorio (\emph{desktop}).} Es el \textbf{fondo grande} que ves al encender. Como una mesa de trabajo: ahí puedes poner íconos de acceso rápido a los programas o archivos que usas a menudo. \emph{Función:} darte un lugar central donde dejar lo que necesitas a la mano. \textbf{Elemento 2 --- La barra de tareas (\emph{taskbar}).} La \textbf{tira horizontal abajo} (en Mac está arriba). Muestra las apps que tienes abiertas en este momento y te deja saltar entre ellas. \emph{Función:} ser el ``índice'' de los programas vivos. \textbf{Elemento 3 --- El menú inicio (\emph{start menu}).} El \textbf{botón Windows abajo-izquierda}. Te muestra todas las apps instaladas, el botón de apagar, búsqueda. \emph{Función:} ser la \emph{puerta principal} para encontrar cualquier cosa en tu computador.}
\milcpaso{2}{milcTurquesa}{\textbf{Elemento 4 --- Las ventanas.} Cuando abres un programa, aparece en un \textbf{rectángulo con barra arriba} que puedes mover, redimensionar, minimizar (esconder), maximizar (pantalla completa) o cerrar (X roja). \emph{Función:} permitir tener varios programas a la vez sin que se mezclen. \textbf{Elemento 5 --- La papelera de reciclaje.} El \textbf{ícono de papelera} en el escritorio. Cuando borras un archivo, va ahí primero (no se borra para siempre). Si te arrepientes, lo recuperas. Si vacías la papelera, ahí sí se pierde. \emph{Función:} darte una segunda oportunidad ante errores de borrado.}
\milcpaso{3}{milcTurquesa}{\textbf{Las 5 funciones del sistema operativo (lo que hace el gerente).} (1) \textbf{Coordina el hardware}: hace que el teclado, ratón, monitor y altavoces hablen con los programas. (2) \textbf{Gestiona archivos}: organiza dónde se guardan tus documentos y los encuentra cuando los necesitas. (3) \textbf{Ejecuta programas}: arranca, mantiene y cierra Word, Chrome, Roblox. (4) \textbf{Maneja usuarios}: si en tu casa varios usan el mismo PC, el SO separa los archivos de cada uno por cuenta. (5) \textbf{Controla la memoria}: decide cuánta RAM le da a cada programa abierto (más a Chrome, menos a la calculadora, por ejemplo).}
\milcpaso{4}{milcTurquesa}{\textbf{Los 5 SO más importantes que vas a encontrar.} (1) \textbf{Windows} (Microsoft): el más común en colegios y casas, fácil de usar. (2) \textbf{macOS} (Apple): elegante, caro, viene en MacBook. (3) \textbf{Linux} (gratuito): muchas versiones, usado en programadores y servidores; \emph{Ubuntu} es la más amigable. (4) \textbf{Android} (Google): el SO de la mayoría de celulares en Colombia. (5) \textbf{iOS} (Apple): el SO del iPhone. \emph{Pista:} no hay un SO ``mejor''. Hay el que te enseñaron a usar y el que pruebas nuevo. Cambiar de SO es como cambiar de barrio: al principio te confundes, después te acomodas.}

\begin{drawbox}{milcTurquesa}{5 elementos del escritorio + 5 funciones del SO}
\textbf{Los 5 elementos visuales:} \\[1mm]
escritorio · barra de tareas · menú inicio · ventanas · papelera. \\[1mm]
\textbf{Las 5 funciones invisibles:} \\[1mm]
coordina hardware · gestiona archivos · ejecuta programas · maneja usuarios · controla memoria.
\end{drawbox}

\begin{softbox}{milcMostaza}{milcGris}{Errores comunes y su corrección}
\textbf{Confusiones que se repiten sobre el SO.} (1) \emph{``Windows es lo mismo que Microsoft Word''} --- Falso. Windows es el sistema operativo (gerente); Word es un programa (un trabajador). (2) \emph{``Si no tengo Windows, no puedo usar Word''} --- Hoy es falso: Word tiene versiones para Windows, Mac, Android, iOS y web. (3) \emph{``El SO se descarga gratis''} --- Depende. Windows cuesta dinero (licencia); Linux es gratis; macOS solo viene con Mac. (4) \emph{``Cuando borro un archivo, se va para siempre''} --- Falso si tiene papelera de reciclaje. Se borra para siempre solo cuando vacías la papelera. (5) \emph{``El SO no se actualiza, está siempre igual''} --- Falso. Los SO se actualizan cada cierto tiempo (Windows Update, actualizaciones de Android) para arreglar errores y agregar mejoras.
\end{softbox}

\begin{infoband}{milcTurquesa}


\textbf{Lo que va al cuaderno (Actividad 2 · EXPLICA).} Título: «Actividad 2 --- Los 5 elementos y las 5 funciones del SO». \textbf{Formato:} 2 tablas pequeñas. Tabla A (5 filas): nombre del elemento + función. Tabla B (5 filas): nombre de la función del SO + ejemplo concreto. \textbf{Extensión:} 10 filas en total. \textbf{Sabes que terminaste cuando} podrías explicar a un primo qué pasa en su computador cuando lo enciende.
\end{infoband}

\puente{Con todo claro, haces el tour anotado y das tu opinión sobre el SO que usas.}

\milcestacion{milcMagenta}{Estación 3 de 3 · Hacer}

\begin{softbox}{milcMagenta}{milcGris}{Cómo vas a construir tu producto}
\textbf{Actividad 3 · APLICA --- Tour anotado del SO + opinión personal} (30 min · individual o en parejas). Vas a hacer 3 cosas en una hoja completa del cuaderno: \textbf{(a)} dibujo del escritorio de Windows con los 5 elementos etiquetados; \textbf{(b)} tabla de 5 funciones del SO con ejemplos tuyos; \textbf{(c)} opinión escrita sobre el SO que más usas tú.
\end{softbox}

\milcpaso{1}{milcMagenta}{\textbf{Paso 1 --- Dibuja el escritorio (mitad superior).} Pon la hoja vertical. En la mitad superior dibuja un \textbf{rectángulo grande} (que ocupa casi toda la mitad) representando la pantalla del computador. Adentro dibuja: un \emph{fondo} con algún detalle (un sol, una montaña, un patrón --- como un wallpaper), \emph{3 íconos} de programas en el escritorio (Word, Chrome, una carpeta), \emph{una barra horizontal abajo} (barra de tareas) con 4 íconos más, \emph{un botón Windows} a la izquierda abajo, y un \emph{ícono de papelera} arriba a la derecha.}
\milcpaso{2}{milcMagenta}{\textbf{Paso 2 --- Etiquetas con flechas.} Desde fuera del rectángulo, traza \textbf{5 flechas} apuntando a los 5 elementos del escritorio. Al final de cada flecha escribe el nombre técnico: \emph{``1. Escritorio (desktop)''}, \emph{``2. Íconos''}, \emph{``3. Barra de tareas''}, \emph{``4. Menú inicio (Windows)''}, \emph{``5. Papelera''}. Si quieres agregar una ventana abierta de un programa, también está bien (sería el elemento extra).}
\milcpaso{3}{milcMagenta}{\textbf{Paso 3 --- Tabla de las 5 funciones (mitad inferior izquierda).} En la mitad inferior izquierda dibuja una tabla de \textbf{5 filas y 2 columnas}. Columna 1: nombre de la función. Columna 2: ejemplo concreto tuyo. Por ejemplo: \emph{``Coordina hardware --- Cuando muevo el ratón, el cursor se mueve en la pantalla.''}, \emph{``Gestiona archivos --- Cuando guardo el Word, va a la carpeta Documentos.''}, etc. Llena las 5 filas.}
\milcpaso{4}{milcMagenta}{\textbf{Paso 4 --- Opinión personal (mitad inferior derecha).} En la mitad inferior derecha escribe un \textbf{párrafo de 5-6 líneas} respondiendo: \emph{``¿Cuál SO uso más yo? (Windows / Android / iOS / Linux / otro). ¿Qué me gusta? ¿Qué me cuesta? ¿He probado otro? Si pudiera probar uno nuevo, ¿cuál sería y por qué?''}. Sé honesto. Esta opinión es tuya, no hay respuesta correcta.}
\milcpaso{5}{milcMagenta}{\textbf{Paso 5 --- Tour real (si tienes computador disponible).} Si en clase tienes acceso a un computador, abre Windows y verifica que sí están los 5 elementos donde dibujaste. Mueve una ventana, abrir y cerrar el menú inicio, abrir la papelera. Esto te confirma lo aprendido. Si no tienes computador hoy, hazlo en casa esta semana.}
\milcpaso{6}{milcMagenta}{\textbf{Paso 6 --- Firma y cierre.} En la esquina inferior derecha firma con nombre completo y fecha. Escribe debajo: \emph{``Ahora sé qué hace el gerente invisible de mi computador.''}. Subraya esa frase.}

\begin{drawbox}{milcMagenta}{Antes de mostrar tu trabajo}
\checkbox La mitad superior tiene el dibujo del escritorio con sus 5 elementos etiquetados. \\[1mm]
\checkbox La tabla de 5 funciones tiene un ejemplo concreto por cada una. \\[1mm]
\checkbox La opinión personal tiene 5-6 líneas y responde las 4 preguntas. \\[1mm]
\checkbox Está firmado con nombre y fecha. \\[1mm]
\checkbox La frase de cierre está subrayada. \\[1mm]
\checkbox Si tuve computador, verifiqué en vivo los 5 elementos.
\end{drawbox}

\begin{softbox}{milcMagenta}{milcGris}{Plantilla de trabajo}
\textbf{Estructura.} \textit{Mitad superior:} dibujo del escritorio + 5 etiquetas con flechas. \textit{Mitad inferior izquierda:} tabla 5 × 2 (función + ejemplo). \textit{Mitad inferior derecha:} párrafo de opinión personal (5-6 líneas). \textit{Esquina inferior:} firma + frase subrayada.
\end{softbox}

\begin{infoband}{milcMagenta}


\textbf{Lo que va al cuaderno (Actividad 3 · APLICA).} Título: «Actividad 3 --- Tour del sistema operativo». \textbf{Formato:} dibujo del escritorio + tabla de funciones + párrafo de opinión + firma. \textbf{Extensión:} 1 hoja entera. \textbf{Sabes que terminaste cuando} podrías mostrar la hoja a un familiar y él entendería qué hace el sistema operativo en el computador.
\end{infoband}

\puente{El producto es el dibujo del escritorio + tabla + opinión, firmado.}

\milcestacion{milcMostaza}{Producto de la sesión}
\textbf{Tour anotado del SO con 5 elementos + 5 funciones + opinión personal · firmado}.
\begin{softbox}{milcMostaza}{milcGris}{Criterios de calidad}
(1) El dibujo del escritorio tiene los 5 elementos etiquetados con flechas. (2) La tabla de 5 funciones tiene ejemplos tuyos. (3) La opinión personal responde las 4 preguntas con honestidad. (4) El tour está firmado con nombre y fecha. (5) La frase de cierre está visible.
\end{softbox}

\puente{Antes de salir, miras tu hoja: ¿podría un primo entender qué es un SO solo leyéndola?}

\milcestacion{milcVino}{Evaluación liberadora · cinco dimensiones}

\begin{softbox}{milcVino}{milcGris}{Sentido de la evaluación}
Evaluar no es solo revisar una nota. Es reconocer qué aprendí, qué cambió en mí, qué emociones manejé, cómo me relaciono con la ciudad y la comunidad, y qué le dejaré a quien viene detrás.
\end{softbox}

\textbf{Mi reflexión liberadora en cinco dimensiones.} Completa cada frase iniciada:

\small
\begin{tabularx}{\textwidth}{>{\bfseries\RaggedRight\arraybackslash}p{3.4cm}Y>{\RaggedRight\arraybackslash}p{4.6cm}}
\rowcolor{milcVino}\color{white}Dimensión & \color{white}\bfseries Mi respuesta (frame starter) & \color{white}\bfseries Algo que descubrí\\
1. Personal & Lo que más cambió en mí fue \linead{6.5cm} & \linead{4.2cm}\\
2. Emocional & La emoción más fuerte fue \linead{2.5cm} y la regulé \linead{3cm} & \linead{4.2cm}\\
3. Ciudadana & En democracia esto importa porque \linead{6cm} & \linead{4.2cm}\\
4. Local (Cartago) & En mi barrio sería útil para \linead{6.5cm} & \linead{4.2cm}\\
5. Intergeneracional & A mi abuela/abuelo le diría \linead{6.5cm} & \linead{4.2cm}\\
\end{tabularx}
\normalsize

\begin{infoband}{milcVino}
\textbf{Para la próxima sustentación.} Vuelve a esta tabla en seis meses. Verás que las respuestas cambian — no porque lo aprendido cambie, sino porque tú cambias. La evaluación liberadora es una conversación contigo a través del tiempo.
\end{infoband}


\puente{Cerramos con 3 voces que te ayudan a entender por qué saber quién coordina la máquina te hace usuario libre.}

\milcestacion{milcGradeDark}{Triángulo de pensamiento · cierre filosófico}

\begin{softbox}{milcVino}{milcGris}{Dussel · lente del nosotros}
\textit{Cuando entiendes quién coordina, entiendes quién manda. Cuando entiendes quién manda, decides si lo aceptas.}

El sistema operativo \emph{decide} cómo se ve tu pantalla, qué apps puedes usar, cómo se organizan los archivos. Cuando lo entiendes, te das cuenta de que \textbf{Windows no es la única forma}: hay Linux (gratis, abierto, donde tú decides), hay otros mundos. Conocer al gerente te permite preguntar: \emph{``¿qué gerente quiero que coordine mi máquina?''}. Esa pregunta es ejercicio de libertad.

\textbf{¿Uso el SO que me asignaron porque me lo asignaron, o ya he probado otro para decidir cuál prefiero?}
\end{softbox}
\linead{\linewidth}\\[1mm]
\linead{\linewidth}

\begin{softbox}{milcMostaza}{milcGris}{Estoicismo · Marco Aurelio · cuidado interior}
\textit{Lo invisible es lo más importante. Lo que no se ve sostiene lo que se ve.}

El sistema operativo es \textbf{invisible la mayor parte del tiempo}. No le ponemos atención porque funciona en silencio. Solo lo notamos cuando falla. \emph{Esto es como la salud, el aire, la familia que está atrás}: tan importantes que no se ven hasta que faltan. Aprender a notar lo invisible te entrena en gratitud y atención. El SO te lo enseña en pequeño.

\textbf{¿A qué cosas invisibles de mi vida les he dejado de poner atención porque ``siempre funcionan''?}
\end{softbox}
\linead{\linewidth}\\[1mm]
\linead{\linewidth}

\begin{softbox}{milcTurquesa}{milcGris}{Floridi · lente de la infoesfera}
\textit{Los sistemas operativos no son neutros: cada uno propone un modo de habitar el mundo digital.}

Cuando usas Windows, vives en un mundo que Microsoft ha diseñado: ese fondo, esa barra, esa lógica de archivos. Cuando usas Android, vives en otro mundo (de Google). Cuando usas iOS, en otro (de Apple). \textbf{Cada SO trae una filosofía}: cómo organizar las cosas, qué te ofrece de regalo y qué te cobra. Aprender a leer esas filosofías te hace ciudadano digital con criterio.

\textbf{¿Qué filosofía me está \emph{enseñando} el SO que más uso, sin que yo lo note?}
\end{softbox}
\linead{\linewidth}\\[1mm]
\linead{\linewidth}

\begin{infoband}{milcNegro}
\textbf{Nota para el docente.} Las tres formulaciones del triángulo son la síntesis del autor de la guía, no citas textuales de los pensadores. Si vas a citarlos en clase, remite a la obra original.
\end{infoband}

\milcestacion{milcVino}{Mi autoevaluación · marca una casilla por fila}

{\small
\setlength{\extrarowheight}{2.2pt}
\begin{tabularx}{\textwidth}{>{\RaggedRight\arraybackslash\bfseries}p{3.0cm}YYY}
\rowcolor{milcVino}\color{white}\bfseries Criterio & \color{white}\bfseries Lo logré & \color{white}\bfseries En proceso & \color{white}\bfseries Necesito apoyo\\[.6mm]
Comprensión del tema & \milcopcion{Explico las ideas con mis palabras.} & \milcopcion{Reconozco las ideas, falta precisión.} & \milcopcion{Necesito repasar lo básico.}\\[.8mm]
\rowcolor{milcGris}Pensamiento computacional & \milcopcion{Usé los pilares al armar mi producto.} & \milcopcion{Usé algunos pilares.} & \milcopcion{Necesito releer la estación 2.}\\[.8mm]
Aplicación y producto & \milcopcion{Mi producto cumple los criterios.} & \milcopcion{Está armado, con ajustes pendientes.} & \milcopcion{Aún está en construcción.}\\[.8mm]
\rowcolor{milcGris}Reflexión 5 dimensiones & \milcopcion{Respondí cada una con honestidad.} & \milcopcion{Respondí algunas con detalle.} & \milcopcion{Mis respuestas son muy generales.}\\[.8mm]
Triángulo de pensamiento & \milcopcion{Conecté las 3 voces con mi trabajo.} & \milcopcion{Conecté al menos una.} & \milcopcion{No las relaciono con mi proceso.}\\[.8mm]
\end{tabularx}}

\vspace{3mm}
\textbf{Hoy entendí que:}\\[1mm]
\linead{\linewidth}\\[1mm]
\linead{\linewidth}

\textbf{Mi compromiso para la próxima sesión es:}\\[1mm]
\linead{\linewidth}\\[1mm]
\linead{\linewidth}

\begin{softbox}{milcMostaza}{milcGris}{Cita de despedida · Marco Aurelio}
\textit{``El obstáculo en el camino se vuelve el camino. Lo que estorba al camino se vuelve camino.''}\\[1mm]
Cada error de hoy, bien observado, se convierte en pieza del camino que pisarás mañana.
\end{softbox}

\small
\begin{tabularx}{\textwidth}{>{\bfseries}p{3.6cm}Y}
\rowcolor{milcGradeDark!12}Nombre completo: & \linead{10cm}\\
Fecha: & \linead{4cm} \quad Curso/grupo: \linead{4cm}\\
Mi firma: & \linead{9cm}\\
\end{tabularx}
\normalsize

\begin{infoband}{milcGradeDark}
\textbf{Para la próxima sesión.} Esta semana, cuando uses tu computador (o el celular), observa: ¿qué está haciendo el SO en cada momento? ¿Cuándo notas su trabajo? La próxima clase aprenderás a \textbf{organizar archivos y carpetas}: el lado del SO que más vas a usar. Si dominas eso, dejarás de perder tus tareas en el desorden y serás más rápido en cualquier computador.
\end{infoband}

\end{document}
